% !TeX root = ../../../main.tex
\subsubsection{Modeling Approaches}
% \vspace{-.25em}
To evaluate performance on the event forecasting and survival prediction tasks, we implement five modeling paradigms: (i) decoder-only large language models (LLMs) with fine-tuned heads, (ii) prompted LLMs without gradient updates (zero- or few-shot), (iii) encoder-only models with task-specific fine-tuned heads, (iv) encoder-masking models with fine-tuning, and (v) encoder-masking models in zero-shot settings.
In Appendix~\ref{ch07:apd:encoder-forecasting} and~\ref{ch07:apd:encoder-survival}, we provide technical details of encoder-based forecasting methods and the survival modeling framework respectively.

\ChSevenParagraph{(i) Fine-tuned LLMs.}  
We apply instruction-tuned decoder-only models from the \textit{Llama} and \textit{DeepSeek} families: \textit{Llama-3.3-70B-Instruct}, \textit{Llama-3.1-8B-Instruct}, and their distilled variants \citep{grattafiori2024llama,liu2024deepseek}.
We also include open-source models with documented training corpora (\textit{OLMO-32B-Instruct} \citep{olmo20242}, \textit{RedPajama-INCITE-7B-Instruct} \citep{weber2024redpajama}) to address potential data leakage concerns---notably, RedPajama explicitly excludes PubMed from its training data.
Additionally, we evaluate a medically fine-tuned decoder model (\textit{MediPhi-PubMed} \citep{corbeil2025modular}) to assess domain-specific benefits.

Each model is paired with a lightweight multilayer perceptron (MLP) head trained for classification or ranking, depending on the task. 
The input consists of a text-formatted prefix of events (e.g., a clinical timeline up to time $t$), optionally accompanied by an instruction template. 
The output layer produces task-specific predictions: binary labels for event occurrence or a permutation over $k$ events for ordering. 
Training is conducted using cross-entropy loss for classification and pairwise ranking loss for ordering.

\ChSevenParagraph{(ii) Prompted LLMs.}  
In the zero- or few-shot setting, we use the same LLM architectures as above, but without any fine-tuning. 
Instead, we supply structured prompts at inference time to guide the model toward task-specific outputs. 
Each prompt includes: 
(1) a system instruction establishing the model’s role (e.g., "You are an expert physician."), 
(2) a user instruction describing the prediction task, and 
(3) one or more input-output examples in a few-shot format.
Prompt designs are customized for each task and are detailed in Appendix~\ref{ch07:apd:prompts}.
Model outputs are parsed to extract binary labels or ordered lists, and errors in generation or ambiguity are excluded from evaluation.

\ChSevenParagraph{(iii) Fine-tuned Encoder-Only Models.}  
We also evaluate a range of encoder-based models trained end-to-end on each task. 
Architectures include general-purpose models (\textit{BERT-base-uncased} \citep{devlin2019bert}, \textit{RoBERTa-base} \citep{liu2019roberta}, \textit{DeBERTa-v3-small} \citep{he2021debertav3}, \textit{ModernBERT-base/large} \citep{warner2024smarter}) and biomedically-pretrained variants (\textit{BioClinical-ModernBERT-base/large} \citep{sounack2025bioclinical}).
For each model, we append a task-specific MLP head for event occurrence or ordering prediction. 
The model input is the same flattened prefix of the event sequence used in other settings, tokenized and formatted according to the respective architecture. 
These models are trained using standard supervised learning objectives and evaluated on the same metrics as other methods: F1 score for event occurrence, and pairwise concordance for temporal ordering. 
For survival analysis, the final hidden states are passed into a time-to-event prediction head and evaluated using the time-dependent concordance index \citep{ch07:antolini2005time}.

\ChSevenParagraph{(iv, v) Encoder-Masking Models.}
We adapt masked language modeling (MLM) for our temporal reasoning tasks.
For event occurrence, models predict masked tokens (``\texttt{yes}''/``\texttt{no}'') in prompts such as ``Will [event] happen within 24 hours? \texttt{[MASK]}''.
For temporal ordering, models predict ``\texttt{before}''/``\texttt{after}'' tokens in prompts comparing event pairs.
We evaluate both fine-tuned variants (trained on task-specific objectives) and zero-shot variants (using pretrained MLM capabilities without additional training).
This approach leverages the bidirectional context of encoder models while maintaining interpretable predictions through constrained vocabulary.

% \vspace{-.5em}
